\chapter{Trabalhos Relacionados}

A classificação automática de arritmias cardíacas tem sido amplamente investigada na literatura recente, com destaque para a transição de métodos 
baseados em extração manual de características para abordagens de aprendizado profundo (\textit{Deep Learning}). Esta seção apresenta estudos 
relevantes que utilizaram Redes Neurais Convolucionais (CNNs) unidimensionais e bidimensionais.

\section{Abordagens Unidimensionais (1D)}

O trabalho de \citeonline{Real-Time} é considerado um dos marcos na aplicação de CNNs para classificação de ECG. Os autores propuseram uma arquitetura de 
CNN1D adaptativa e específica para o paciente, capaz de extrair características diretamente do sinal bruto, sem a necessidade de pré-processamento 
complexo. O estudo demonstrou que redes 1D podem superar métodos tradicionais em velocidade e eficiência, alcançando acurácia superior a 97\% na 
classificação de batimentos ventriculares e supraventriculares. Este trabalho fundamenta a escolha da CNN1D como uma das abordagens avaliadas nesta 
pesquisa.

Seguindo essa linha, \citeonline{DeepTransferableRepresentation}. propuseram uma representação profunda transferível, utilizando redes residuais aplicadas a sinais de 1D. Esse 
trabalho propõe reutilizar os aprendizados obtidos do treinamento para detecção de arritmias para detectar infarto do miocárdio. 
A acurácia média obtida foi de 93.4\% para a detecção de arritmias. Um ponto a ser levado em conta é a complexidade da rede utilizada, que teve 13 camadas.
Reutilizando o aprendizado, foi obtido uma acurácia de 95.9\% para detecção de infarto do miocárdio.

Em uma abordagem mais recente, \citeonline{DeepResnet} desenvolveram um modelo de rede neural residual profunda convolucional unidimensional 
(Deep ResNet 1D) para classificar cinco classes de batimentos cardíacos da AAMI. Utilizando a base de dados PhysioNet MIT-BIH Arrhythmia, os autores 
lidaram com o desbalanceamento aos dados aplicando a técnica de sobreamostragem sintética da classe minoritária (SMOTE) no conjunto de 
treinamento. Avaliado através de validação cruzada de 10 vias (\textit{10-fold cross-validation}), o modelo alcançou uma acurácia média de 98,63\%, 
evidenciando a robustez das redes residuais unidimensionais para a extração direta de características de sinais de ECG.

Um ponto de discussão metodológica central que tem ganhado força na literatura recente sobre a base MIT-BIH é a distinção entre os protocolos 
intra-paciente e inter-paciente. Trabalhos intra-paciente costumam apresentar resultados superestimados devido ao vazamento de dados, já que batimentos 
de um mesmo indivíduo podem aparecer simultaneamente no treino e no teste. Visando superar essa limitação com uma avaliação inter-paciente rigorosa, 
\citeonline{Zhou2023} propuseram uma abordagem combinando um Perceptron de Múltiplas Camadas (MLP), uma rede de Cápsulas de Peso e um modelo 
Sequência-para-Sequência aplicados ao sinal 1D. O estudo alcançou uma acurácia de 99,88\% e se destacou por lidar de forma eficaz com o desbalanceamento 
das classes minoritárias da base de dados sem recorrer a técnicas de aumento artificial, como o SMOTE.

\section{Abordagens Bidimensionais (2D)}

Com o objetivo de explorar informações no domínio da frequência, diversas pesquisas investigaram a transformação de sinais 1D em representações 2D. 
\citeonline{UsingSTFT-BasedSpectrogram} desenvolveram um método baseado na Transformada de Fourier de Curta Duração (STFT) para converter batimentos de ECG em espectrogramas, 
que serviram de entrada para uma CNN 2D. Os autores relataram uma acurácia média de 99,00\%, argumentando que a representação tempo-frequência 
facilita a distinção de arritmias com características espectrais complexas.

Diversificando as técnicas de transformação tempo-frequência, \citeonline{wavelettransformwithdeeplearningmodel} exploraram o uso da Transformada Wavelet Contínua (CWT) para converter os 
sinais de ECG em escalogramas, que serviram de entrada para uma CNN 2D. Diferentemente da STFT utilizada nesta monografia, a CWT oferece uma resolução 
multiescala, o que permitiu aos autores atingir uma acurácia global de 99,65\%. Apesar do excelente desempenho, a geração de escalogramas via CWT 
é matematicamente mais custosa que a STFT, o que pode impactar aplicações em tempo real.

Avançando nas metodologias de classificação, arquiteturas como os Transformers têm sido adaptadas para o processamento de imagens bidimensionais no 
domínio médico. \citeonline{Chen2024} introduziram um método combinando mapas de tempo-frequência, gerados através da Transformada Wavelet Contínua 
com a ondaleta complexa de Morlet, com o modelo Swin Transformer. Esta arquitetura hierárquica utiliza um mecanismo de autoatenção baseado em 
janelas deslocadas para capturar de forma eficiente tanto as características locais quanto as dependências globais de longo alcance presentes nas 
imagens do ECG. Avaliado na base MIT-BIH, o modelo alcançou uma acurácia de 99,34\% no cenário intra-paciente e de 98,37\% numa avaliação rigorosa 
inter-paciente. Este estudo evidencia o grande potencial dos Transformers visuais na superação das limitações de dependência temporal das redes 
tradicionais.

Buscando unir a capacidade de extração de padrões espaciais das CNNs com a modelagem de dependências temporais, \citeonline{Franklin2022} apresentaram 
uma arquitetura híbrida CNN-LSTM aplicada a representações 2D. O método converte os sinais de ECG em imagens em escala de cinza, as quais são processadas 
por uma CNN 2D para identificar assinaturas morfológicas, seguidas por camadas LSTM que capturam o contexto sequencial e as dependências de longo prazo 
dos batimentos. Esta abordagem demonstrou ser eficaz na redução do impacto de ruídos biológicos, permitindo uma classificação robusta.

\section{Comparações}

A Tabela \ref{tab:trabalhos_relacionados} resume o comparativo entre os métodos do estado da arte.

\begin{table}[ht]
    \caption{Comparativo entre trabalhos relacionados}
    \label{tab:trabalhos_relacionados}
    \centering
    \small
    \resizebox{\textwidth}{!}{
    \begin{tabular}{ccccc}
        \hline
        \textbf{Autor/Ano} & \textbf{Entrada} & \textbf{Arquitetura} & \textbf{Classes} & \textbf{Acurácia Média} \\
        \hline
        Kiranyaz et al. (2016) & Sinal 1D & CNN 1D (Rasa) & 5 & $\approx$ 97.40\% \\
        Kachuee et al. (2018) & Sinal 1D & ResNet 1D & 5 & 93.40\% \\
        Khan et al. (2023) & Sinal 1D & Deep ResNet 1D & 5 & 98.63\% \\
        Zhou et al. (2023) & Sinal 1D & MLP + WCapsule + Seq2Seq & 5 & 99.88\% \\
        Huang et al. (2019) & STFT & CNN 2D & 5 & 99.00\% \\
        Mohonta et al. (2022) & CWT & CNN2D & 5 & 99.65\% \\
        Chen et al. (2024) & CWT & Swin Transformer & 5 & 98.37\% (Inter), 99,34\% (Intra) \\
        Franklin et al. (2022) & Imagem 2D & CNN-LSTM & 5 & 99.33\% \\
        \hline
    \end{tabular}
    }
    \source
\end{table}